\begin{proposition}[3.3.6]
\label{IV.2.3.3.6-fr}
Sean $f:X\to S$, $g:Y\to S$ dos morfismos, $\mathcal F$ un
$\mathcal O_X$-módulo cuasicoherente y $\mathcal G$ un $\mathcal O_Y$-módulo
cuasicoherente. Se supone que: 1\textsuperscript{o} $\mathcal G$ es plano
respecto de $g$; 2\textsuperscript{o} $X$ es localmente noetheriano y, para todo
$s\in S$, $g^{-1}(s)$ es localmente noetheriano (lo que ocurrirá si $Y$ es
también localmente noetheriano). Sea $Z=X\times_S Y$; para todo par $(x,y)$
tal que $x\in X$, $y\in Y$ y

\oldpage[IV-2]{45}
$f(x)=g(y)=s$, sea $T_{x,y}$ el preesquema
$\operatorname{Spec}(\boldsymbol{k}(x)\otimes_{\mathbf{k}(s)}\boldsymbol{k}(y))$, y sea $I_{x,y}$ la imagen de
$\operatorname{Ass}(\mathcal O_{T_{x,y}})$ por el monomorfismo canónico
$T_{x,y}\to Z$ \textup{(\textbf I, 3.4.9)}. Se tiene entonces
\end{proposition}

\begin{env}[3.3.6.1]
\label{IV.2.3.3.6.1-fr}
\[
\operatorname{Ass}(\mathcal F\otimes_S\mathcal G)
=\bigcup_{x\in\operatorname{Ass}(\mathcal F)}
\left(\bigcup_{y\in\operatorname{Ass}(\mathcal G_{f(x)})}I_{x,y}\right)
\]
\end{env}

\textit{donde, para todo} $s\in S$, $\mathcal G_s=\mathcal G\otimes_{\mathcal O_S}\boldsymbol{k}(s)$.

Sean $p:Z\to X$, $q:Z\to Y$ las proyecciones canónicas, de modo que se tiene
el diagrama conmutativo
\[
\xymatrix{
X \ar[d]_f & Z \ar[l]_p \ar[d]^q \\
S & Y \ar[l]^g
}
\]
Pongamos $\mathcal G'=q^*(\mathcal G)$, de modo que
$\mathcal F\otimes_S\mathcal G=\mathcal F\otimes_X\mathcal G'$; como
$\mathcal G'$ es plano respecto de $p$ (2.1.4), se deduce de (3.3.1) que

\begin{env}[3.3.6.2]
\label{IV.2.3.3.6.2-fr}
\[
\operatorname{Ass}(\mathcal F\otimes_S\mathcal G)
=\bigcup_{x\in\operatorname{Ass}(\mathcal F)}
\operatorname{Ass}(\mathcal G'_x)
\]
\end{env}

con $\mathcal G'_x=\mathcal G'\otimes_{\mathcal O_X}\boldsymbol{k}(x)$. Si $s=f(x)$, se tiene
$\mathcal G'_x=\mathcal G_s\otimes_{\mathbf{k}(s)}\boldsymbol{k}(x)$, y
$p^{-1}(x)=g^{-1}(s)\otimes_{\mathbf{k}(s)}\boldsymbol{k}(x)$; además, como el cuerpo $\boldsymbol{k}(x)$ es un
$\boldsymbol{k}(s)$-módulo plano, el morfismo $p^{-1}(x)\to g^{-1}(s)$ es plano (2.1.4);
aplicando (3.3.1) a este morfismo, se obtiene

\begin{env}[3.3.6.3]
\label{IV.2.3.3.6.3-fr}
\[
\operatorname{Ass}(\mathcal G'_x)
=\bigcup_{y\in\operatorname{Ass}(\mathcal G_s)}
\operatorname{Ass}(\mathcal O_{T_{x,y}})
\]
\end{env}

de donde se sigue la proposición.
